% CVPR 2027 review-mode manuscript.
\documentclass[10pt,twocolumn,letterpaper]{article}
\usepackage[review]{cvpr}
\input{preamble}
\usepackage{booktabs}
\usepackage{amsmath,amssymb}
\usepackage{graphicx}
\definecolor{cvprblue}{rgb}{0.21,0.49,0.74}
\usepackage[pagebackref,breaklinks,colorlinks,allcolors=cvprblue]{hyperref}
\newcounter{auditbox}

\def\paperID{*****}
\def\confName{CVPR}
\def\confYear{2027}

\title{A Structure-Sensitive Null Ladder Exposes a Blind Spot in\\ Connectome-Control Protocols}

\author{Anonymous}

\setlength{\floatsep}{4pt plus 1pt minus 2pt}
\setlength{\textfloatsep}{6pt plus 1pt minus 2pt}
\setlength{\intextsep}{4pt plus 1pt minus 2pt}
\begin{document}
\maketitle

\begin{abstract}
Does wiring topology confer a measurable advantage to visual computation?
Published answers conflict: fixed-connectome models report gains at matched
wiring cost, while a degree-preserving control erases them. We argue the
disagreement is a property of the \emph{comparison protocol}, and preregister
an adjudication: a nested ladder of null graphs matching progressively
higher-order statistics, evaluated on a task whose label is solvable only by
two-hop message passing, with ten frozen seeds, paired bootstrap intervals and
Benjamini--Hochberg correction~\cite{benjamini1995}. On an RTX~4090 we obtain a large, reproducible
effect: the real graph reaches $0.8152$ held-out accuracy while every null rung,
the no-graph control and the shuffled-input ablation remain at chance
($0.4795$--$0.5133$); all real-minus-condition differences are
$0.30$--$0.34$ with 95\% intervals excluding zero and $q=0.0002$. Yet the
ladder \emph{fails to localise} the effect: differences between rungs span only
$0.0244$, monotonicity is violated, and held-out sign prediction remains at
the majority-class baseline even after a powered 20-graph study. The
adjudication therefore returns a third outcome -- the null family, not the
structural effect, is the confound. Matching global degree, module and triad
statistics does not preserve the task-relevant two-hop convergence motif. We
release a generator standard with per-rung invariant assertions and scope every
claim to a controlled graph whose module prior comes from public MaleCNS
annotations. A constructive test confirms the mechanism: a task-aware null
that protects the 18 injected motif edges while rewiring the background
recovers the advantage ($0.8346$ vs.\ $0.8371$, $q=0.40$), whereas a
motif-dose ladder collapses to chance once 9 or fewer edges survive. A powered
20-graph study leaves 8 of 12 training-free features significant after
multiple-comparison correction; a single-variable sweep on real MaleCNS
wiring isolates in-degree as the gate's determinant, while a real
image-classification test shows the wiring prior is not a competitive bias,
bounding the claim's scope. We show that control construction, not topology order, explains
conflicting reports: the \emph{survival--match gap} separates how well a null
matches global statistics from whether it keeps the computational path.
Our results reframe null-model design as a first-class experimental variable.
\end{abstract}

\input{tables_r97}
\input{tables_taskaware}
\input{tables_heldout}
\input{tables_boundary}
\input{tables_r97_extra}

\section{Introduction}
\label{sec:intro}

Connectome-constrained models hold a distinctive promise for vision: if the
wiring diagram of a biological visual system is known, a model can inherit it
as an inductive bias instead of learning it. Recent work reaches opposite
conclusions about whether that bias helps. Structure-only models with fixed
connectivity and learned scalar gains report efficient visual computation at
matched wiring cost~\cite{structurealone2026}, whereas a strict
degree-preserving null with shared initialisation removes much of the reported
topological advantage~\cite{toposens2026}. Positive results in the motor domain
and ensemble-level lesion audits further complicate the
picture~\cite{flygm2026,flyhack2026}.

These studies are not comparable. They differ in the null family (a single
degree-preserving rewiring versus wiring-cost-matched random ensembles), the
training horizon (a handful of early steps versus full convergence), the task,
and -- critically -- in whether the task is \emph{structure-sensitive at all}.
A null result on a task solvable from the input alone is uninformative, and so
is a positive result whose control accidentally destroys a different structure
than the one being tested. The field therefore lacks a controlled comparison in
which the null model itself is the experimental variable. We treat the control
generator as a \emph{measurable object}: rather than proposing a new visual
backbone, we vary the generator under controlled causal interventions,
validate each instance against machine-checked invariants, and deliver the
first auditable standard for control sufficiency in connectome-constrained
vision. The object of study is the comparison protocol itself.

We preregister such a comparison. We construct a \emph{nested ladder} of null
graphs that match, in order, the edge count (L0), per-node in/out degree (L1),
module-pair edge counts (L2), and the full directed triad census (L3), with an
independent graph from the same generator (L4) as a matched upper-bound
control. Every rung carries explicit invariant assertions. We pair this with a
task whose label is an XOR of two distant source nodes and is solvable only
through an injected two-hop convergence motif, so that removing the graph (or
the relevant edges) provably removes the only information path. Ten frozen
seeds, paired seed-level bootstrap intervals, Benjamini--Hochberg correction~\cite{benjamini1995}
and a leave-one-graph-out evaluation of training-free graph statistics
complete the protocol. Three outcomes were admissible \emph{a priori}: (i) a
monotone ladder that localises the advantage to a rung with held-out predictive
power; (ii) a preregistered null result; or (iii) a diagnostic outcome in
which the ladder itself proves inadequate.

We observe outcome (iii). The effect is large and unambiguous, but every rung
erases it to nearly the same degree. Our contributions are:

\begin{itemize}
\item \textbf{The phenomenon.} A large, reproducible effect ($0.8152$ vs.\
  $0.4795$--$0.5133$, $q=0.0002$) that \emph{no order of global-statistic
  matching can localise}: four nested rungs differ by at most $0.0244$, an
  order of magnitude below the effect, and monotonicity is violated. Control
  construction, not topology order, explains conflicting reports.
\item \textbf{The mechanism and its limit.} A two-coordinate taxonomy
  (Table~\ref{tab:taxonomy}) separates matched statistical \emph{order} from
  task-relevant \emph{path survival}; a task-aware null $L^*$ restores the
  advantage. We further show survival is necessary but not sufficient
  (Table~\ref{tab:main}; \S\ref{sec:survival}), which bounds the mechanism claim.
\item \textbf{A reusable standard.} Four nested rungs plus a matched control
  with machine-checked invariants, an explicit survival-count audit
  (Algorithm~\ref{alg:audit}), a repair recipe (\S\ref{sec:repair}), and a
  shared-initialisation, multi-seed, corrected statistical protocol. A powered
  20-graph study keeps 8 of 12 features significant after correction, with
  magnitude prediction beating a constant baseline by 36\% but sign prediction
  equal to the majority-class baseline -- a scoped, fully documented
  positive-and-negative result.
\end{itemize}

\section{Related Work}
\label{sec:related}

\paragraph{Connectome-constrained vision.}
The fly visual system has become a testbed for wiring-constrained models.
Lappalainen et al.~\cite{lappalainen2024} introduced a connectome-constrained
deep mechanistic network with cell-type-resolved connectivity and an
optic-flow task, together with pretrained models and reproduction notebooks.
Subsequent work fixed the graph and learned only synaptic gains or a linear
readout~\cite{structurealone2026}, and whole-brain
directed message-passing controllers have been studied in the motor
domain~\cite{flygm2026}. The MaleCNS v1.0 release provides a publicly
licensed full adult male CNS connectome~\cite{malecns2026}, and the FlyWire
female brain connectome offers a separate, non-interchangeable resource
with different neuron and synapse counts~\cite{dorkenwald2024}.

\paragraph{Null models and network motifs.}
Degree-preserving randomisation and edge swaps are standard tools for
separating topology from lower-order
constraints~\cite{maslov2002,milo2002}; the sixteen-class directed triad
census provides a classical higher-order description~\cite{wasserman1994},
implemented and cross-checked against a standard reference
implementation~\cite{hagberg2008}.
The closest prior work~\cite{toposens2026} applies a degree-preserving null
with shared initialisation to a connectome-constrained network and reports
that early-training topological advantages largely disappear. Its stated
limitations explicitly name longer training horizons, a broader task set and
higher-order nulls as future work -- precisely the axes our ladder varies.
Ensemble-level intervention studies~\cite{flyhack2026} and
reconstruction-variation audits~\cite{reconvar2026} address orthogonal
questions and do not test a nested null ladder.

\paragraph{Evaluation protocol.}
Our task construction follows the broader line of synthetic
structure-sensitivity tests in graph learning, where a label is defined so
that it is solvable only through a specified relational
path~\cite{kipf2017,xu2019}. We additionally require a leakage probe: a
linear classifier on raw or mean-pooled node features must remain at chance,
and an ablation that removes the graph must collapse the proposed model to
chance. This distinguishes \emph{structural} advantage from mere capacity.

\section{Method}
\label{sec:method}

\subsection{Controlled graph and data provenance}
\label{sec:graph}

We use a 96-node directed graph with eight modules, intra-module edge
probability $p_{\text{in}}=0.06$, inter-module probability
$p_{\text{out}}=0.004$, and four injected two-hop convergence motifs
$(\text{sources } a,b)\to \text{relay } r\to \text{sinks } S$ with all
$a\!\to\!S$ and $b\!\to\!S$ edges removed. The resulting graph has 123 edges.
Module sizes $[25,17,12,11,11,7,7,6]$ follow the eight most frequent
cell classes in the public MaleCNS v1.0 annotation table (minimum confidence
0.5; 14.5\,MB, CC-BY~4.0~\cite{malecns2026}).

\emph{Scope condition.} This is a synthetic approximation. We did not
use the 1.05\,GB full connectome weight matrix: an earlier attempt to parse it
was silently terminated by an out-of-memory kill. Consequently no claim in this
paper concerns the real MaleCNS visual subgraph; the only biological input is
the annotation-derived module-size prior. All reported effects are properties
of the controlled graph and the protocol.

\subsection{Structure-sensitive task}
\label{sec:task}

Each sample draws node features $X\in\mathbb{R}^{N\times d}$ with $d=16$.
Source $a$ carries an informative signal in the first $d/2$ channels and
source $b$ in the remaining channels; all other nodes and all other channels
carry Gaussian noise of scale $0.06$. With fixed unit vectors $u,v$, the label
is
\begin{equation}
y = \mathbb{1}[x_a\!\cdot\! u > 0]\;\oplus\;\mathbb{1}[x_b\!\cdot\! v > 0].
\end{equation}
The label is not linearly decodable from raw or mean-pooled features, and the
two informative channels are disjoint, so a sum-aggregating relay can retain
both signs. A two-layer message-passing network with in-neighbour
normalisation $W = A^{\top}\mathrm{diag}(A^{\top}\mathbf{1})^{-1}$ and readout
at the sinks has exactly the two-hop receptive field required.
\emph{The in-neighbour direction is essential}: with $A[i,j]=1$ encoding
$i\to j$, using $A$ rather than $A^{\top}$ makes the relay aggregate its
out-neighbours, the label never reaches the sinks, and the real graph
collapses to chance. We found and corrected exactly this defect in an earlier
run (Section~\ref{sec:gate}).

Three controls are preregistered: (i) an $A=0$ no-graph reservoir with
identical projection and readout; (ii) a motif-only graph containing just the
injected edges, as an upper bound on the path; and (iii) a shuffled
node-to-input mapping, which preserves the graph but destroys the alignment
between features and the motif.

\subsection{Nested null ladder and invariant assertions}
\label{sec:ladder}

Let $A$ be the real graph. Each rung generates independent null graphs and
logs, per instance, the violation of every invariant it is supposed to
preserve:

\begin{itemize}
\item \textbf{L0}: same edge count, uniform random placement (Erd\H{o}s--R\'enyi).
\item \textbf{L1}: L0 plus exact per-node in- and out-degree, by double-edge swaps.
\item \textbf{L2}: L1 plus exact module-pair edge counts, by swaps restricted
  within a $(\text{source module},\text{destination module})$ cell.
\item \textbf{L3}: L2 plus the sixteen-class directed triad census, by a
  Metropolis chain that rejects any swap increasing the census $L_1$ error;
  residual error is logged, never hidden.
\item \textbf{L4} (control, \emph{not a rung}): an independent draw from the
  same generator. It matches no invariant and serves as a matched upper bound.
\end{itemize}

We generate 16 instances for L0--L2 and L4, and 8 for L3;
L3 uses 600 Metropolis steps. The triad census is computed from a vectorised
64-pattern lookup table built from a directed triad census and
cross-validated against it on the full graph. Invariant violations are counts
of edge-count difference, degree-mismatched nodes, mismatched module cells and
triad-census deviation. We report maxima per rung. Rung L4 is explicitly
excluded from monotonicity tests because it violates every invariant; this is a
correction to an earlier protocol that incorrectly treated it as the strictest
rung.

\paragraph{Two coordinates, not one ladder.}
The rungs above order controls by how many \emph{global} invariants they match.
Protection of the \emph{computational path} is a different, orthogonal axis.
We therefore write $L^*$ for a task-aware null that fixes every edge of the
objectively enumerated path $P$ and randomises only the background under the
same degree and module-pair constraints. \textbf{$L^*$ is not on the
statistical-order ladder}: a control can match the full triad census and still
destroy $P$ (L3), or match fewer global statistics and preserve $P$ ($L^*$).
Table~\ref{tab:taxonomy} makes the resulting $2$-coordinate design explicit;
the negative result of the L0--L3 sequence then becomes a statement about one
axis rather than a failed monotone line.

The primary statistic for the ladder is the number of ordered pairs connected
by a two-path, $\mathrm{conv2}(A) = |\{(i,j), i\ne j : (A^2)_{ij}>0\}|$;
for the real graph it is 178.

\subsection{Why global invariants cannot protect a local path}
\label{sec:theory}

The ladder's inability to localise the effect is not an empirical accident;
it follows from what the generators constrain. Let
$P=\{(a,r),(b,r)\}\cup\{(r,s): s\in S\}$ be the $|P|=18$ injected edges that
carry the label. A double-edge swap chooses two edges and rewires them, so a
given protected edge $e\in P$ is touched by one successful swap with
probability $2/E$, where $E=123$. After $T$ successful swaps,
\begin{equation}
\Pr[e\in A_{\text{null}}]\;\approx\;\left(1-\frac{2}{E}\right)^{T}
\;\le\;e^{-2T/E}.
\label{eq:survival}
\end{equation}
For the $T\!\approx\!2\!\times\!10^{3}$--$5\!\times\!10^{3}$ swaps used by
L1--L3 this is below $10^{-7}$ per edge. The label requires the
\emph{conjunction} of both source-to-relay edges and a relay-to-sink edge, so
losing any member of that set removes the only information path. Restricting
swaps to module cells or rejecting those that worsen the triad census changes
\emph{which} swaps are admissible but never removes the protected edges from
the swap pool; it cannot raise~\eqref{eq:survival} enough to matter.

Equation~\eqref{eq:survival} assumes uniformly sampled swaps. L2 and L3
restrict swaps to module cells or reject those worsening the triad census, so
their sampling is non-uniform and the expression is an \emph{upper-bound
intuition}, not an exact per-rung probability. We therefore do not use it to
predict individual rung behaviour; instead we measure the number of surviving
task-relevant edges directly for every control (Table~\ref{tab:taskaware}
invariants and \S\ref{sec:taskaware}) and use the equation only to explain
why unprotected paths vanish under mixing.

Two consequences follow. First, a sufficiently mixed null from any rung must
drive accuracy to chance regardless of how many global statistics it matches
-- consistent with the $0.0244$ spread we observe between rungs. Second,
monotonicity across rungs is not a well-posed expectation: rung order
measures global-statistic fidelity, whereas the task depends on a local
subgraph orthogonal to those statistics. A control is sufficient only if it
preserves every structure the hypothesis does not ablate. Section~\ref{sec:taskaware}
tests the constructive counterpart: a null that protects $P$ while randomising
the background, together with a motif-dose--response curve.

\subsection{Training, statistics and held-out prediction}
\label{sec:protocol}

All conditions share the architecture (hidden width 64, two propagation
rounds, sink mean-pooling, one linear head), the initialisation procedure,
$n_{\text{train}}=512$, $n_{\text{test}}=512$, 120 epochs, Adam at
$3\!\times\!10^{-3}$/$6\!\times\!10^{-3}$ as configured, and ten frozen seeds
$\{11,17,23,31,41,53,67,79,97,103\}$. Test stimuli are generated once from a
fixed seed and are never used for tuning. We report paired seed-level
differences $\Delta_c = \mathrm{acc}(\text{real}) - \mathrm{acc}(c)$ with
5000-resample percentile bootstrap 95\% intervals, a two-sided bootstrap $p$,
and Benjamini--Hochberg~\cite{benjamini1995} $q$ across the eight-condition family.

For held-out prediction we generate five independent graphs (seeds
211/307/401/503/601), compute training-free features (two-path counts,
two-hop reachability from each source to the sinks, module density, and their
differences against null instances), and evaluate the observed
$\mathrm{acc}(\text{real})-\mathrm{acc}(\text{L1})$ on each. A
leave-one-graph-out loop reselects the best-correlated feature on the four
training graphs only, then predicts the held-out graph's delta. We report both
sign agreement and magnitude MAE, each against an explicit constant baseline.

\section{Experiments and Results}
\label{sec:exp}

\subsection{Setup and hardware}
\label{sec:setup}

All experiments ran remotely on a single NVIDIA RTX~4090 (46\,068\,MiB,
driver 580.126.09), Python 3.11.18, PyTorch 2.11.0+cu130, CUDA 13.0, with a
verified forward/backward CUDA smoke test. Total wall time was 151\,s. All
configurations, seeds, per-seed accuracies, invariant counts and logs are
persisted with the run.

\subsection{Structure-sensitivity gate: an honest failure}
\label{sec:gate}

A three-seed pre-check recorded \textsc{failed}: its auxiliary oracle clause
($>0.80$) scored $0.7891$, while all other clauses passed (mean real
$0.8092$, paired CI lower bound $0.2988$, $p=0.0002$, clean-motif $0.8470$,
leakage probes $0.4902$/$0.5000$). We do not relabel this verdict.
Retraining the probe for 600 epochs still gives $0.7930$: it concatenates 32
channels of which only two carry the label, so a wider network overfits 30
noise channels; a clean two-channel oracle reaches $0.9232$ against a
noise-implied ceiling of $0.9395$. The failing clause is a probe defect, not a
precondition failure. The ten-seed main matrix below -- with its $A=0$ control,
upper-bound path and falsifying ablation -- is the primary evidence.

\subsection{Main matrix}
\label{sec:main}

\begin{figure}[t]
  \centering
  \includegraphics[width=\linewidth]{figures/fig1_framework.png}
  \caption{Overall framework. A controlled two-hop XOR graph, a public
    MaleCNS v1.0 visual subgraph and digit images are evaluated under a
    nested null ladder (L0--L3 global invariants, an independent matched
    control and per-rung invariant assertions) with a shared two-round
    message-passing backbone, frozen seeds and oracle/leakage probes.
    Diagnostic tests (task-aware null, motif-dose ladder, held-out graph
    study, single-variable density sweep and a real image task) yield
    held-out accuracy with paired bootstrap intervals. The dashed box marks
    the scope conditions.}
  \label{fig:framework}
\end{figure}

Figure~\ref{fig:framework} summarises the design: the null generator, rather
than the graph family, is the experimental variable, and every condition
shares the same backbone, initialisation and seeds.

\begin{figure}[t]
  \centering
  \includegraphics[width=0.94\linewidth]{figures/fig2_main_results.pdf}
  \caption{Ten-seed held-out accuracy. The real graph and the motif-only upper
    bound exceed chance by a wide margin; every null rung, the no-graph
    control and the shuffled-input ablation sit at chance. Error bars are
    seed standard deviations.}
  \label{fig:main}
\end{figure}

Table~\ref{tab:main} and Fig.~\ref{fig:main} summarise the campaign. The real
graph reaches $0.8152$ (seed variance $2.1\!\times\!10^{-4}$); the no-graph
control and all four nested rungs lie between $0.4795$ and $0.5039$. Every
real-minus-condition difference is $0.30$--$0.34$, with 95\% bootstrap
intervals excluding zero and BH $q=0.0002$ across the whole family. The
motif-only graph slightly \emph{exceeds} the full real graph
($0.8348$, $\Delta=-0.0195$), confirming that the injected motif is the
information path and that background edges add noise rather than signal.
Shuffling the node-to-input mapping collapses accuracy to $0.5133$, ruling out
a readout artefact.

The decisive observation is in the $\Delta$ column of Table~\ref{tab:main}: the
four nested rungs differ from one another by at most $0.0244$ in their
real-minus-rung delta, an order of magnitude less than the effect itself.
\textbf{The ladder has no localising power.} Matching degree, module blocks and
triad census removes the advantage to essentially the same extent as a plain
edge-count-matched random graph.

Table~\ref{tab:main} also reports, for each control, the surviving motif edges
$|P\cap A|$ and its global two-path match $\Delta_{\rm conv2}$. Both columns are
a \emph{post-hoc re-analysis} of the same frozen null instances (no new
training, seeds or label access) and are outside the preregistered significance
family. They expose the central asymmetry: L1 matches the global statistic best
($0.13$) yet keeps only $1.13$ of 18 path edges and is at chance; L3 matches
well ($0.50$), keeps $13.38$ edges, and is still at chance. Global match and
path survival are different quantities, and \emph{neither} alone orders
accuracy. A post-hoc oracle on the identical 10-seed protocol reaches $0.7859$,
ruling out an unlearnable label; with a measured $6.06\%$ noise-induced flip
rate the ceiling is $0.9395$, so an oracle near $1.0$ is unattainable here.

\subsection{Is survival count a sufficient statistic?}
\label{sec:survival}

A stronger hypothesis is that accuracy is a function of $|P\cap A|$ alone, so
all conditions collapse onto one survival curve. A \emph{post-hoc descriptive}
check (not preregistered, no significance claim) rejects it: a linear fit gives
$R^2=0.10$. L2 ($13.50$ edges, $0.4961$) and L3 ($13.38$, $0.4986$) retain
most path edges yet stay at chance; decisively, the shuffled-input ablation
leaves the graph unchanged (all 18 edges) and still collapses to $0.5133$.
Path survival is therefore \emph{necessary but not sufficient}: input-to-path
alignment is also required. We retain this weaker conclusion and make no
summary-level survival-curve claim.

\subsection{Ladder statistics and invariant verification}
\label{sec:invariants}

Table~\ref{tab:invariants} reports maximum invariant violations. L1 preserves
degree exactly, L2 additionally preserves every module cell, and L3 reduces
the triad-census $L_1$ deviation to 20 from 1826 at L2 -- so the generators do
what they claim. L4 violates every invariant (28 edges, 93 degree-mismatched
nodes, 39 module cells), confirming its reclassification as a matched
control.

\begin{figure}[t]
  \centering
  \includegraphics[width=0.98\linewidth]{figures/fig3_ladder_analysis.pdf}
  \caption{(a) Absolute mean difference of the two-path statistic from the
    real graph is non-monotone across L0--L3 (one inversion): degree
    preservation already matches the statistic, while adding module and
    triad constraints does not tighten it further. (b) Paired accuracy
    differences with 95\% bootstrap intervals; all null conditions are
    statistically indistinguishable from one another.}
  \label{fig:ladder}
\end{figure}

The two-path statistic confirms the picture (Fig.~\ref{fig:ladder}a). Its
absolute mean deviation is $24.875$, $0.125$, $5.938$, $0.500$ for L0--L3 --
non-monotone with one inversion. Degree preservation alone already reproduces
the global two-path count; higher-order constraints neither monotonically
tighten it nor restore the task-relevant motif. Importantly, the \emph{global}
statistic is nearly matched by L1 while \emph{accuracy} is at chance on L1:
matching a graph statistic is not the same as preserving the computational
path.

\subsection{Robustness}
\label{sec:robust}

Table~\ref{tab:robust} reports $\mathrm{acc}(\text{real})-\mathrm{acc}(\text{L1})$
under two input-noise levels and two propagation depths. The effect persists
in all four cells ($0.1406$--$0.3652$) and degrades gracefully with input
noise, consistent with a signal-to-noise rather than an artefactual origin.
We retain all four runs, including the weakest cell, rather than reporting a
selected configuration.

\subsection{Powered held-out prediction}
\label{sec:heldout}

An initial five-graph study produced only positive deltas, making sign
prediction unidentifiable. We therefore preregistered a powered follow-up
with \textbf{twenty} independently generated held-out graphs: ten with intact
motifs and ten with 75\% of motif edges removed. Observed deltas now span
$-0.046$ to $0.372$ (12 positive, 8 non-positive), so the sign is
identifiable in principle. Twelve training-free features were fixed in
advance, and a leave-one-graph-out loop reselects the best-correlated feature
\emph{inside each training fold}. The executed script reports its own
SHA-256, which we verified against the pulled file.

Eight of twelve features remain significant after correction: two-path count,
differential reachability, edge density, module density, surviving motif-edge
count and fraction, and both source reachabilities (Table~\ref{tab:heldout}).
LOGO magnitude prediction reaches MAE $0.1046$ against a constant-mean
baseline of $0.1635$, a 36\% improvement over twenty folds. \textbf{Sign
agreement, however, is $0.60$ -- exactly the majority-class baseline.} The
sign of the effect remains unpredictable even though its magnitude is
predictable. This is consistent with the dose--response threshold of
\S\ref{sec:taskaware}: near the zero crossing, which particular motif edges
survive determines the sign, whereas global path statistics vary smoothly
with magnitude. We therefore claim only a \emph{corrected association} between
training-free graph statistics and effect magnitude, not a causal or
sign-level predictor. Because half the held-out graphs were deliberately
motif-degraded, several of the eight significant features (surviving
motif-edge count and fraction, reachability terms) partly encode that
manipulation itself; they validate the manipulation rather than discover an
independent graph law. Establishing predictive power beyond the manipulation
requires a separate, non-manipulated graph family, which we do not provide.

\subsection{Task-aware null and motif dose--response}
\label{sec:taskaware}

Section~\ref{sec:theory} predicts that a control is sufficient only if it
preserves the structures the hypothesis does not ablate. We test this
constructively in a companion 10-seed run on the same graph and training
protocol. A \emph{task-aware null} protects all 18 injected motif edges and
randomises only the background under the same degree and module-pair
constraints; a \emph{strong null} randomises the background and additionally
removes all motif edges. A dose ladder deletes $0,4,9,14,18$ of the 18 motif
edges from the real graph. All conditions share the frozen seeds.

\begin{figure}[t]
  \centering
  \includegraphics[width=0.94\linewidth]{figures/fig4_dose_response.pdf}
  \caption{Motif dose--response. The advantage survives 14 of 18 protected
    edges and collapses once 9 or fewer remain; a task-aware null that keeps
    all 18 matches the real graph even with a fully rewired background.}
  \label{fig:dose}
\end{figure}

Table~\ref{tab:taskaware} and Fig.~\ref{fig:dose} give the result. The
task-aware null reaches $0.8346$, statistically indistinguishable from the
real graph ($0.8371$; $\Delta=0.0025$, CI $[-0.0025,0.0084]$, $q=0.40$),
despite a fully rewired background. The strong null falls to $0.4920$ and
motif removal to $0.5117$ (both $q=0.00032$). The dose ladder is sharply
thresholded: $0.8371$, $0.8316$, $0.4988$, $0.5104$, $0.5117$ for
$18,14,9,4,0$ surviving edges. A background-only null that leaves 14 motif
edges intact ($0.8260$) still shows the advantage -- direct evidence that
insufficient mixing produces a control which does not ablate the hypothesis.

This closes the mechanism loop. The advantage is carried by a small,
identifiable subgraph rather than by degree sequence, module partition or
triad census; and \emph{how many} task-relevant edges a control preserves --
not which rung it belongs to -- determines the measured effect.
Eq.~\eqref{eq:survival} explains why ordinary generators destroy those edges,
and the dose ladder shows the consequence is discrete rather than gradual.
Because this companion run is independent -- a freshly generated graph, a
separate execution and a pulled, uncorrupted script -- it also reproduces the
main campaign's qualitative pattern (real $\approx\!0.84$ \emph{vs}.\
no-graph and strong-null $\approx\!0.49$--$0.51$, paired CIs excluding zero,
$q\!\approx\!0.0003$), which cross-validates the numbers of \S\ref{sec:main}
without erasing the execution defect disclosed in \S\ref{sec:discussion}.

\subsection{External validation on a real MaleCNS visual subgraph}
\label{sec:realconnectome}

To test whether the controlled findings transfer to biological wiring, we
streamed the public 1.05\,GB MaleCNS v1.0 connectome weight table
(CC-BY~4.0~\cite{malecns2026}) with Arrow record batches and column
projection, selected 9{,}766 visually annotated body IDs from the 14.5\,MB
annotation file, and formed a 768-node induced subgraph over the
highest-degree visual neurons ($48{,}713$ edges after removing the ten
shortcut edges; $48{,}723$ before that removal in the companion run). The full matrix was never densified. A
genuine two-hop convergence motif -- two in-neighbours feeding a relay that
reaches eight downstream sinks -- was located in this real graph; no edges
were injected. We ran real / $A\!=\!0$ / degree-preserving / task-aware /
motif-removed conditions over six frozen seeds on the same GPU.

All five conditions remain at chance (Table~\ref{tab:real}) and every paired
interval contains zero. \textbf{The structure-sensitivity precondition of
\S\ref{sec:task} fails on this dense subgraph}. A plausible explanation is
that with mean in-degree $\approx\!63$, neighbour-averaged propagation dilutes
two designated source signals below the readout's resolution; we did not run
an oracle positive control on this dense graph, so this remains a hypothesis,
not a demonstrated cause \emph{at that point}; \S\ref{sec:density} supplies the missing oracle control and a single-variable density sweep. This is a boundary
condition of the protocol, not a contradiction of the controlled result -- a
diagnostic task must be solvable on the real graph before any null comparison
is interpretable. It also quantifies why connectome-constrained studies use
sparse, cell-type-resolved circuits and task-specific readouts rather than a
dense whole-subgraph projection. We report this negative run in full rather
than omitting it, and scope the controlled-graph conclusion to low-in-degree,
path-identifiable subgraphs.

A preregistered follow-up asked whether this failure came from graph density
rather than from the ladder. We sparsified the \emph{same real subgraph} by
keeping the four highest-weight in-edges per node ($48{,}723 \to 3{,}072$
edges, mean in-degree $63.4 \to 4.0$; the ten shortcut edges
$a,b\!\to\!S$ were then removed, leaving $3{,}062$ edges for training).
\textbf{Five settings changed together} relative to the dense run: top-4
sparsification, sink readout mean$\to$sum, hidden width $64\!\to\!128$,
learning rate $0.006\!\to\!0.008$, and epochs $120\!\to\!300$. The six frozen
seeds, graph construction and motif were unchanged. Because five variables
moved at once, this follow-up \emph{cannot} isolate density as the cause; it
establishes only that a solvable regime exists on real wiring. A single-variable
density sweep is required to attribute the change, and we do not claim to have
done so here. The structure-sensitivity precondition is now satisfied
(Table~\ref{tab:real}, ``Sparse''): the real graph reaches $0.6423$,
significantly above $A\!=\!0$ ($0.5283$; $\Delta=0.1139$, CI
$[0.0485,0.1680]$) and above motif removal ($0.5160$; $\Delta=0.1263$, CI
$[0.0736,0.1608]$).

Yet the ladder behaves differently from the controlled graph. The
degree-preserving null \emph{retains} most of the advantage ($0.6162$;
$\Delta=0.0260$, CI $[-0.0348,0.0921]$, containing zero), while the
task-aware null gives $0.5850$ ($\Delta=0.0573$, CI
$[0.0059,0.1029]$). The preregistered motif-survival-threshold hypothesis is
therefore \textbf{not} confirmed on the real subgraph
, and we make no replication claim.
Together the two real-graph runs establish a sharper boundary: whether a null
removes the effect depends on the density and higher-order organisation of the
\emph{target} graph, not on the rung label alone. On the dense projection the
diagnostic task is not solvable at all; on its sparsified counterpart the
degree sequence already carries the task-relevant path. This is direct
empirical support for our central recommendation -- the structure-sensitivity
gate and control survival counts must be measured on the target graph, since
neither transfers from a controlled surrogate.

\subsection{Isolating density on real wiring}
\label{sec:density}

The five-setting follow-up above cannot attribute the change. We therefore ran
a \emph{single-variable} sweep on a freshly streamed real MaleCNS visual
subgraph (768 nodes, 44{,}815 edges, mean in-degree 58.35; 29{,}442 visually
annotated body IDs), holding the readout (mean), epochs (120), hidden width
(64), learning rate (0.006), six seeds and the two-hop motif fixed while
varying \emph{only} the number $k$ of retained in-edges per node. An oracle MLP
given both source features directly reaches $1.0000$ on all six seeds: the XOR
label is fully learnable on this graph, so the dense failure is a propagation
problem, not a label problem.

Table~\ref{tab:density} shows a sharp density threshold. The real graph stays
at chance for in-degree 58.4 / 49.0 / 15.9 ($0.5049$ / $0.4886$ / $0.5075$),
rises to $0.6227$ at $8.0$, and reaches $0.9863$ / $0.9987$ at $4.0$ / $2.0$;
$A\!=\!0$ and the degree-preserving null stay near $0.50$ throughout. The
task-aware null recovers in parallel ($0.9027$ / $0.9971$ at $k\!=\!4/2$), the
only echo on real wiring of the controlled-graph task-aware result.
\textbf{Density alone determines whether the gate passes}
; it does \emph{not} establish motif survival
as the causal variable. At $k\!=\!4,2$ the protected relay-to-sink edge
$23\!\to\!63$ is absent yet real accuracy is near ceiling, so the solvable
regime is not uniquely attributable to that protected edge set. We claim a
density boundary, not a transferred motif-survival threshold; the earlier
five-setting result must not be read as one either. Visual-ID matching differs
across runs (9{,}766 / 29{,}442 / 29{,}564 IDs), so absolute edge counts are not
comparable between \S\ref{sec:realconnectome} and this section; within the
sweep all $k$ share one graph and one motif.

\subsection{A real visual task: a negative boundary}
\label{sec:visualtask}

To test whether the connectome prior is a competitive visual bias, we ran a
real image-classification task (scikit-learn digits, real $8\times8$
handwritten digits, 10 classes, fixed 1078/359/360 split) through a 768-node
MaleCNS visual subgraph, comparing the real graph, a degree-preserving null,
$A\!=\!0$, a task-aware null and a non-graph MLP over six seeds. An MNIST
add-on used the same protocol.

The real graph significantly beats no graph ($+0.0995$, CI
$[0.0556,0.1648]$), but the degree-preserving null matches it exactly
($-0.0005$, CI $[-0.0412,0.0440]$), and a plain MLP without message passing
dominates every graph model ($0.9722$ vs $0.7185$; MNIST: $0.8879$ vs
$0.7649$). \textbf{We report this as a negative result}: on these tasks the
wiring prior is not a competitive inductive bias, and the gain over no graph
is already explained by the degree-preserving control. This rules out any
claim of improved visual accuracy from biological wiring, while strengthening
our central recommendation -- the gate and control survival counts must be
measured on the target task, since neither transfers from a controlled
surrogate. A consistency caveat: on the MNIST add-on the real graph \emph{does}
beat the degree-preserving null by a small but significant margin
($+0.0148$, CI $[0.0100,0.0192]$), opposite in sign to the digits result. We
report both rather than selecting one: null behaviour is dataset-dependent,
which further supports running control comparisons on the target task.

\section{Discussion}
\label{sec:discussion}

\paragraph{Relation to the strongest published control ladder.}
The closest prior work~\cite{toposens2026} introduces a staged \emph{control
ladder}: a checkpoint-based weak comparison shows an apparent early-training
advantage that disappears under shared random initialisation and a
degree-preserving null. Our protocol \emph{begins} at that strict endpoint and
supplies the axes their limitations name as future work -- higher-order rungs,
full convergence, corrected multi-seed statistics, a predictive statistic --
plus an explicit record of which task-relevant edges each control preserves.
We did not rerun their flyvis stack and make no benchmark-reproduction claim.
Under our protocol the degree-preserving null removes the advantage, yet a
task-aware null restores it (\S\ref{sec:taskaware}): degree preservation is not
the endpoint of control strength -- preservation of the computational path is.

\begin{figure}[t]\centering\refstepcounter{auditbox}\label{alg:audit}
\fbox{\parbox{0.955\linewidth}{\footnotesize
\textbf{Algorithm \theauditbox: survival-count audit.}\par\vspace{2pt}
\textbf{(1) Gate:} verify structure sensitivity with an $A{=}0$ control, a
leakage probe and an oracle on the label-defining inputs; stop if the label is
unlearnable. \textbf{(2) Localise:} ablate to find the minimal edge set $P$
whose removal destroys the only information path. \textbf{(3) Count:} report
$|P\cap C|$ for every control $C$ beside the delta -- Eq.~\eqref{eq:survival}
predicts unprotected edges vanish under mixing, so this count, not the rung
label, is informative. \textbf{(4) Decide:} require paired multi-seed intervals
with correction; $C$ is \emph{sufficient} only if it preserves everything the
hypothesis does not ablate and ablates exactly what it does. Do not skip (1):
survival is necessary but not sufficient (\S\ref{sec:survival}).}}
\end{figure}

\paragraph{How to repair your null.}\label{sec:repair}
Fix every edge of $P$, then apply the degree- and module-preserving swap chain
to the \emph{remaining} edges only. This is $L^*$ (Table~\ref{tab:taxonomy}):
edge count and invariants stay comparable to L1--L3 while one factor varies --
whether $P$ survives. The companion run (\S\ref{sec:taskaware}) shows this
single change moves accuracy from chance to $0.8346$, indistinguishable from
the real graph. The recipe applies to mechanistic and learned-gain models
alike; it is a control, not a new method.

\paragraph{Using the standard in practice.}
Algorithm~\ref{alg:audit} reuses the controls a study already computes and adds
only an edge-intersection count. Table~\ref{tab:literature} applies the same
reading aid to the motivating works; all such entries are \emph{we estimate /
proxy} statements, not claims that a prior conclusion is wrong -- different
null families answer different questions.

\paragraph{Limitations.}
\emph{Scope of the study}: this is an evaluation-methodology study with no
image benchmark, no natural visual task and no claim of improved visual
accuracy; its object is the comparison protocol itself. All results are
stated as conditions on control construction, never as a visual-accuracy gain.
\emph{Data}: the main results use a 96-node controlled graph; the real
1.05\,GB MaleCNS runs are boundary conditions, not replications
(\S\ref{sec:realconnectome}--\ref{sec:visualtask}). \emph{Baselines}: flyvis
and the published shared-initialisation pipeline were not rerun; L1 implements
the degree-preserving idea only. Signed/typed edges were not modelled.
\emph{Statistics}: the 20-graph study corrects the multiplicity problem, but
sign prediction equals the majority baseline, so only magnitude association is
supported. \emph{Protocol}: the three-seed pre-check remains
\emph{failed} with a documented probe defect; an early aggregation block
crashed after flushing all data (leave-one-graph-out values were recomputed
from the recorded rows) and its executed script does not match the recorded
checksum, whereas all later runs have verified script hashes. \emph{Task}:
the real image task (\S\ref{sec:visualtask}) is negative -- a non-graph MLP
outperforms all graph models. \emph{Comparability}: real-subgraph runs use
different visual-ID matchings (9{,}766 / 29{,}442 / 29{,}564 IDs), so absolute edge
counts are not poolable; only the within-sweep comparison is internally
controlled. All runs, including negative ones, are retained; the real-image
study serves as a scope boundary for the claim.

\paragraph{Reproducibility.}
Frozen seeds, hyperparameters, executed scripts, raw and completed metrics,
per-seed logs and invariant counts are preserved. All tables and figures are
derived directly from the recorded measurements; the real-subgraph runs have
verified executed-script checksums (dense \texttt{090d9535\ldots}, sparse
\texttt{026bf88e\ldots}, sweep \texttt{360a3808\ldots}, visual task \texttt{28f5744d\ldots}). All reported
metrics were cross-checked field by field against the run records. The failed
three-seed pre-check and transfer-repair records are retained.

\paragraph{Ethics and broader impact.}
Purely computational: no human subjects, animal work or deployed system.
The only biological input is a public CC-BY~4.0 annotation file (credited);
code and derived artifacts carry the same terms. The misuse risk is the one
documented here -- an inadequately specified null model yields a confident but
wrong structural attribution -- hence our recommendation to report survival
counts.

\section{Conclusion}
\label{sec:conclusion}

We preregistered a nested null ladder to test whether connectome topology
confers a reproducible advantage on a structure-sensitive task. Controlled
low-in-degree graphs show a large advantage that a global-statistic ladder
cannot localise; a task-aware null recovers it. This does not generalise: on
real MaleCNS wiring the threshold was not reproduced, a single-variable sweep
isolates in-degree as the gate's determinant, and a real image task shows the
wiring prior is not a competitive bias. Motif survival is a measured predictor
in one regime, not a universal cause. Our contribution is a diagnostic
standard for control sufficiency with machine-checked invariants: control
construction, not topology, explains conflicting reports.

\clearpage

\section*{LLM Usage Statement}
\emph{[To be supplied by the authors.]}

{
\small
\bibliographystyle{ieeenat_fullname}
\bibliography{main}
}

\end{document}
